% !PRINTED 506
\weberpage{506}
%<<CONT
\[
a_1E_1\Phi_1 = P_1, \quad a_2E_2\Phi_2 = P_2
\]
become primitive functions, and then
\weq{2}{a_1a_2E_1E_2\Phi = P_1P_2}
is also a primitive function. From (1) and (2) it follows that
\weq{3}{a_1a_2E_1E_2P = aEP_1P_2}
and consequently
\weq{4}{a = a_1a_2.}
If therefore $a=1$, the natural numbers $a_1, a_2$ must also be $=1$, and we have the theorem:

4. \gesperrt{If $\omega$ is an integral functional in $\Omega$ and $\Phi(t)$ a function of $t$ with integral coefficients in $\overline{R}$ which vanishes for $t=\omega$, then every rational divisor of $\Phi(t)$ also has integral coefficients in $\overline{R}$; in particular, the irreducible function $F(t)$, of which $\omega$ is a root according to §.\,136, 8., has integral rational functionals as coefficients.}

This theorem is of fundamental importance because it teaches that a functional $\omega$ which in the field $\Omega$ belongs among the fractional ones cannot, as might be supposed, appear in a higher field as an integral functional.

If a functional $\omega$ is not integral, it satisfies an equation
\[
\Phi(\omega) = \omega^m + A_1\omega^{m-1} + A_2\omega^{m-2} + \cdots + A_m = 0,
\]
in which the coefficients $A_1, A_2, \ldots, A_m$ are rational, but not integral, functionals. If $a$ is the common denominator of the absolute values of $A_1, A_2, \ldots, A_m$, then $aA_1, aA_2, \ldots, aA_m$ are integral functionals. Now
\[
a^m\Phi(\omega) = (a\omega)^m + aA_1(a\omega)^{m-1} + a^2A_2(a\omega)^{m-2} + \cdots
+ a^mA_m = 0,
\]
and $a\omega$ is therefore an integral functional. From this it follows:

5. \gesperrt{Every functional $\omega$ of the field $\Omega$ can be transformed into an integral functional by multiplication with an integral rational number, and therefore every functional $\omega$ can be represented as a quotient of two integral functionals. This representation is possible in infinitely many different ways, among others in such a way that the denominator is a natural number.}
